\chapter{Giới thiệu}

Mạng di động đang chuyển từ phần cứng chuyên dụng sang phần mềm trên máy chủ thông thường.
Các chức năng từng chạy trên thiết bị tích hợp theo chiều dọc, chuyên biệt cố định, nay được
tái hiện thực thành chương trình trên máy tính đa mục đích. Mạng truy nhập vô tuyến (RAN),
lâu nay là phần phụ thuộc phần cứng nhất trong ngăn xếp, cũng nằm trong xu hướng đó.

Trong 5G, trạm gốc, tên chuẩn là next-generation NodeB (gNB), được tách thành một đơn vị
trung tâm và một hoặc nhiều đơn vị phân tán. Đơn vị trung tâm lại được tách thành các chức
năng mặt điều khiển và mặt người dùng~\cite{3gpp38401,3gpp38300}. Open-RAN đẩy cùng ý tưởng
xa hơn: RAN trở thành phần mềm ảo hoá, đa nhà cung cấp, nối với nhau qua giao diện
mở~\cite{polese2023oran}. Mặt điều khiển thu được, mặt điều khiển đơn vị trung tâm gNB
(gNB-CU-CP), là một chương trình lớn, có trạng thái, tương tranh cao. Nó theo dõi mọi thiết bị
đang kết nối, mỗi thiết bị mà các chuẩn gọi là thiết bị người dùng (UE), chạy báo hiệu Điều
khiển tài nguyên vô tuyến (Radio Resource Control, RRC)~\cite{3gpp38331}, và điều phối di
động cùng chuyển giao, tất cả trong giới hạn độ trễ nghiêm ngặt.


Tải tại một site bận rộn rất cao. Một site có thể chạy hàng nghìn thủ tục quản lý kết nối và
di động mỗi giây, mỗi thủ tục phải hoàn thành trong ngân sách cỡ mili giây mà hệ thống 5G đặt
cho báo hiệu mặt điều khiển~\cite{3gpp38913}. Một mặt điều khiển triển khai có thể theo dõi
hàng triệu thiết bị hoạt động cùng lúc. Khi RAN chuyển sang triển khai đám mây, tách rời, tính
đúng đắn, độ tin cậy và thời gian của phần mềm này ngày càng quyết định chất lượng của toàn
mạng. Một sự cố hoặc một đỉnh độ trễ đuôi ở đây không phải lỗi cục bộ. Mọi thiết bị site phục
vụ đều bị ảnh hưởng~\cite{dean2013tail,barroso2017microseconds}.


Phần mềm mặt điều khiển này, hiện nay, chủ yếu được viết bằng C++. Lựa chọn đó là hợp lý:
C++ mang lại kiểm soát mức thấp, công cụ trưởng thành và hiệu năng mà hệ thống nhạy độ trễ
đòi hỏi. Nhưng hiệu năng đó đi kèm chi phí đã được biết. C++ đặt gánh nặng quản lý bộ nhớ và
tương tranh lên lập trình viên, và hàng thập kỷ kinh nghiệm cho thấy gánh nặng này thường
xuyên bị xử lý sai, sinh ra lỗi use-after-free, double-free, tranh chấp dữ liệu và các lỗi
liên quan, tốn kém để tìm, nguy hiểm khi vận hành, và chiếm phần lớn lỗ hổng bảo mật trong
thực tiễn công nghiệp~\cite{msrc2019memory,google2022android}. Rust đã nổi lên như ngôn ngữ
hệ thống đại chúng đầu tiên hứa hẹn thoát khỏi đánh đổi này: an toàn bộ nhớ và không tranh
chấp dữ liệu được thực thi lúc biên dịch trên nền tảng đã được máy
kiểm~\cite{klabnik2023rust,jung2018rustbelt}, với hiệu năng dự kiến cạnh tranh với
C++~\cite{balasubramanian2017rust}. Liệu lời hứa đó có đúng với mặt điều khiển cấp viễn thông,
nhạy độ trễ, nơi khả năng dự đoán khi quá tải quan trọng ngang tốc độ thuần, vẫn chưa được
hiểu rõ. Luận văn này, thực hiện tại Nokia, trả lời câu hỏi trực tiếp: phát triển hai bản hiện
thực tương đương về chức năng của một phân hệ RAN đại diện, một dịch vụ Quản lý Ngữ cảnh UE
mô hình hoá theo một thành phần gNB-CU-CP thật của Nokia, bằng C++ hiện đại và Rust, rồi so
sánh chúng dưới khối lượng báo hiệu tổng hợp sát thực tế.


Ba thuật ngữ lặp lại xuyên suốt. \emph{gNB-CU-CP} là phần mềm mặt điều khiển mà công trình này
xoay quanh. \emph{Quản lý Ngữ cảnh UE (UECM)} là phân hệ cụ thể được tái hiện thực và đo
lường. \emph{Độ trễ đuôi}, các phân vị p99, p99.9 và p99.99 chứ không phải trung bình, là thuộc
tính quyết định, hơn cả tốc độ trung bình, liệu hệ thống như vậy có đáp ứng hạn
chót~\cite{dean2013tail,barroso2017microseconds}.



\section{Động lực}\label{sec:motivation}
Ba xu hướng khiến phép so sánh này đúng thời điểm. Xu hướng thứ nhất là tiến hoá đám mây của
RAN. Các chức năng mặt điều khiển nay chạy trên máy chủ đa mục đích, nên phần mềm sống lâu
hơn, tương tranh cao hơn, và phải đáp ứng hạn chót mà không còn tính xác định của phần cứng
chuyên dụng. Tách rời Open-RAN biến trạm gốc thành phần mềm ảo hoá, đa nhà
cung cấp~\cite{polese2023oran}. Một hồi quy độ trễ hoặc một lỗi hỏng bộ nhớ khi đó là rủi ro
toàn mạng, và lựa chọn ngôn ngữ là một quyết định kỹ thuật.

Xu hướng thứ hai là trí tuệ nhân tạo (AI), vốn từng được mạng dùng để tự tinh chỉnh và nay
được mạng mang như một khối lượng công việc riêng~\cite{polese2026airan}. International
Mobile Telecommunications-2030 (IMT-2030), khung 6G của Liên minh Viễn thông quốc tế
(International Telecommunication Union, ITU), liệt kê vai trò thứ hai đó như một kịch bản sử
dụng~\cite{itu2023imt2030}. Kiến trúc hình thành quanh nó là AI-RAN, trong đó công việc AI và
vô tuyến dùng chung một máy chủ gia tốc~\cite{airan2026orchestrator}. Mặt điều khiển ở đó
không còn chiếm trọn máy chủ. Nó chia sẻ bộ xử lý trung tâm (CPU), bộ nhớ và cache với công
việc suy luận mà nó không lập lịch. Khi hàng xóm đó chiếm dung lượng, mặt điều khiển bão hoà
sớm hơn: nó ngừng theo kịp nhu cầu ở tải thấp hơn so với khi chạy một mình.
Chương~\ref{chapter:eval} tìm khoảng cách rộng nhất giữa hai ngôn ngữ sau điểm đó.

Xu hướng thứ ba là chuyển sang ngôn ngữ an toàn bộ nhớ. Các tổ chức phần mềm lớn báo cáo rằng
phần lớn lỗ hổng bảo mật nghiêm trọng bắt nguồn từ lỗi an toàn bộ nhớ, và việc áp dụng ngôn
ngữ an toàn bộ nhớ làm giảm chúng~\cite{msrc2019memory,google2022android,vanderstoep2025android,google2024securebydesign}.
Công trình độc lập lập luận rằng an toàn bộ nhớ nay cần yêu cầu được chuẩn hoá nếu muốn được
mua sắm~\cite{watson2025memorysafety}. Các cơ quan an ninh mạng kêu gọi rời C và C++, và Đạo
luật Cyber Resilience của Liên minh châu Âu (EU) thêm nghĩa vụ pháp lý từ tháng 9 năm 2026 và
tháng 12 năm 2027~\cite{cisa2023memorysafe,oncd2024building,eu2024cra}. Đạo luật không nêu tên
ngôn ngữ lập trình nào. Bằng chứng trực tiếp cho một mặt điều khiển RAN vẫn còn thiếu. Cung
cấp bằng chứng đó là động lực của công trình này.

\section{Bối cảnh công nghiệp}\label{sec:industrial}
Luận văn này được thực hiện tại Nokia. Công ty thuộc số ít nhà cung cấp xây, duy trì và phát
triển phần mềm RAN sản xuất từ đầu đến cuối. Năm nhà cung cấp lớn nhất chiếm 96\% doanh thu
RAN toàn cầu trong ba quý đầu năm 2025, với mức tập trung theo chỉ số Herfindahl--Hirschman
cao nhất trong mười năm~\cite{delloro2025ran3q,delloro2026ran}. Phân hệ được nghiên cứu, UECM,
được mô hình hoá theo một thành phần gNB-CU-CP thực chứ không phải ví dụ sách giáo khoa.

Trong kiến trúc của 3rd Generation Partnership Project (3GPP), gNB-CU-CP là nửa mặt điều khiển
của đơn vị trung tâm. Nó lưu RRC và quản lý vòng đời ngữ cảnh
UE~\cite{3gpp38401,3gpp38331}, dưới ngân sách cỡ mili giây mà hệ thống 5G đặt cho báo hiệu điều
khiển~\cite{3gpp38913}. Chương~\ref{chapter:background} đặt nó trong kiến trúc rộng hơn. Vì
được viết bằng C++, nó thừa hưởng hiệu năng của ngôn ngữ đó cùng các mối nguy an toàn bộ nhớ
và tương tranh. Ba xu hướng của Mục~\ref{sec:motivation} gặp nhau ở đây. Phần mềm lớp này chạy
trên mạng nhà khai thác đang vận hành. Lựa chọn ngôn ngữ là mối quan tâm sản xuất.

UECM đủ nhỏ để tái hiện thực trung thành và lý giải trọn vẹn. Nó vẫn mang những thuộc tính làm
mã mặt điều khiển khó. Hiện thực bám theo một đặc tả công khai, đã cố định, không phải mã độc
quyền, nên độc giả ngoài Nokia có thể tái lập phép so sánh.

\paragraph{Một quyết định có tầm một thế hệ.}
Một lựa chọn ngôn ngữ trong phần mềm lớp này sống lâu hơn thế hệ mà nó được đưa ra. 3GPP đã
dành Release 20 cho các nghiên cứu 6G và Release 21 cho công việc chuẩn tắc. ITU yêu cầu đề
xuất công nghệ vào đầu 2029 và nộp đặc tả vào giữa 2030~\cite{3gpp2024road6g}. Công việc đó
bao gồm di trú từ 5G và liên vận hành với 5G. Một mặt điều khiển viết hôm nay được duy trì
cùng thế hệ kế, chứ không bị loại. IMT-2030 vẫn nêu an ninh và khả năng phục hồi như nguyên
tắc thiết kế và vẫn mang kịch bản độ trễ thấp siêu tin cậy ở dạng được tăng
cường~\cite{itu2023imt2030}.

Các phép đo ở đây dùng một phân hệ thế hệ hiện tại và các yêu cầu đã công bố. Không tuyên bố
nào được đưa ra về những yêu cầu chưa được viết.

\paragraph{Một dạng lỗi cụ thể.}
Hãy xét một cơn bão di động, nhiều UE chuyển giao cùng lúc. Một ngữ cảnh UE được giải phóng
khi nhả, trong khi một thông điệp báo hiệu đang bay vẫn còn tham chiếu tới nó, là lỗi
use-after-free. Ở C++ điều đó biên dịch và có thể vượt các kiểm thử thông thường. Nó rồi có
thể làm hỏng trạng thái của một UE khác hoặc làm sập luồng xử lý dưới tương tranh cao. Cùng
lỗi đó có thể hiện ra như một đỉnh độ trễ đuôi thay vì một sự cố. Ở Rust an toàn, thiết kế
không biên dịch: sở hữu và mượn từ chối tham chiếu treo. Đó là lý do luận văn này đo cả
\emph{an toàn} (những lỗi nào được ngăn, và khi nào) lẫn \emph{khả năng dự đoán} (độ trễ đuôi
khi quá tải), chứ không chỉ thông lượng trung bình.

\section{Phát biểu bài toán}
Hai thuộc tính của miền này làm lựa chọn ngôn ngữ mang hệ quả. Thứ nhất, quản lý bộ nhớ và
tương tranh thủ công trong C++ phơi phần mềm trước các lỗi như use-after-free, double-free,
tranh chấp dữ liệu và vô hiệu hoá iterator, im lặng lúc biên dịch và tốn kém để tìm và sửa ở
quy mô viễn thông. Thứ hai, cái quyết định liệu hạn chót RAN có được đáp ứng không phải thông
lượng trung bình mà là hành vi trường hợp xấu nhất và đuôi dưới tải và quá
tải~\cite{dean2013tail}, một thuộc tính dễ hồi quy và khó đo đúng. Rust hứa loại lớp lỗi thứ
nhất tại thời điểm biên dịch trong khi tuyên bố hiệu năng cạnh tranh trên lớp thứ hai, nhưng
tuyên bố đó chủ yếu dựa trên các benchmark đa mục đích chứ không trên khối lượng viễn thông
có trạng thái, tương tranh, nhạy độ trễ.

Gộp các mạch đó lại, bài toán cụ thể mà luận văn giải quyết như sau. Các nhà cung cấp viễn
thông đang chịu áp lực ngày càng lớn, vừa kỹ thuật vừa ngày càng mang tính điều tiết, phải áp
dụng ngôn ngữ an toàn bộ nhớ, nhưng với mặt điều khiển RAN họ thiếu bằng chứng \emph{trực tiếp,
có kiểm soát} để dựa vào quyết định đó. Chưa được xác lập liệu một hiện thực Rust của một phân
hệ mặt điều khiển có trạng thái, tương tranh có thể khớp C++ hiện đại về \emph{độ trễ đuôi và
khả năng dự đoán khi quá tải} (khác với thông lượng trung bình), liệu các bảo đảm lúc biên
dịch của Rust có thực sự loại những lớp lỗi cụ thể tác động tới mã C++ sẵn có, và chi phí áp
dụng tăng dần \emph{tốn} bao nhiêu bên trong một mã nguồn C++ lớn. Bằng chứng sẵn có phần lớn
là đa mục đích và trường hợp trung
bình~\cite{balasubramanian2017rust,qin2020rust}, để lại ba câu hỏi này (khả năng dự đoán, phạm
vi lỗi được ngăn, và chi phí di trú) còn mở cho một phân hệ đại diện. Trả lời chúng trong điều
kiện có kiểm soát, tái lập được là bài toán luận văn này đặt ra để giải.

\section{Câu hỏi nghiên cứu và giả thuyết}\label{sec:rqs}
Luận văn được tổ chức quanh bốn câu hỏi nghiên cứu. Chúng đặt khả năng dự đoán độ trễ lên
trước, rồi phạm vi lỗi, rồi chi phí tài nguyên, rồi áp dụng tăng dần:

\begin{description}
  \item[RQ1] Rust và C++ hiện đại so sánh thế nào về khả năng dự đoán độ trễ (trung vị, độ trễ
        đuôi p99, p99.9 và p99.99, và cực đại) cho một khối lượng mặt điều khiển RAN dưới tải
        tăng dần và quá tải?
  \item[RQ2] Những lớp lỗi an toàn bộ nhớ và tương tranh xảy ra trong C++ nào bị loại tại thời
        điểm biên dịch trong hiện thực Rust, và các lớp còn lại được phát hiện ở giai đoạn nào
        với chi phí nào?
  \item[RQ3] Chi phí, nếu có, của mô hình an toàn Rust (kiểm biên, bản sao do sở hữu, chi phí
        trừu tượng) trên thông lượng, mức dùng CPU, và hành vi bộ nhớ là bao nhiêu?
  \item[RQ4] Những hệ quả thực tiễn (công sức, khả năng tương tác qua giao diện hàm ngoại
        (FFI), và rủi ro) của việc áp dụng Rust tăng dần trong một hệ sinh thái phần mềm RAN
        C++ sẵn có là gì?
\end{description}

Ba giả thuyết được viết ra trước mọi phép đo.
\begin{description}
  \item[H1] Với khối lượng này hai ngôn ngữ tương đương về hiệu năng dưới bão hoà, mọi khác
        biệt nếu có chỉ nằm ở đuôi xa.
  \item[H2] Rust an toàn ngăn tại thời điểm biên dịch những lớp lỗi an toàn bộ nhớ và tương
        tranh mà C++ chỉ cố bắt sau khi chương trình đã được xây.
  \item[H3] Thời gian thêm khi C++ gọi sang Rust nhỏ so với chi phí của chính các thao tác, nên
        chi phí thực của việc thêm Rust từng phần là công việc kỹ thuật, không phải tốc độ lúc
        chạy.
\end{description}
Chương~\ref{chapter:discussion} xem lại cả ba dựa trên bằng chứng.

Danh mục lỗi cho RQ2 gồm các lớp mà Rust an toàn được kỳ vọng ngăn và các lớp mà nó được kỳ
vọng không ngăn. Nghiên cứu khi đó có thể tìm nơi bảo đảm an toàn dừng lại, thay vì chỉ cho
thấy bảo đảm đó giữ.

Một nghiên cứu trước đã ước lượng biên đó từ một danh mục, không từ các chương trình được biên
dịch. Nó lấy Common Weakness Enumeration (CWE), một danh sách dài các kiểu lỗi, và đánh dấu
mỗi kiểu là kiểu mà Rust an toàn sẽ ngăn hoặc sẽ không. Nó rồi đếm các lỗ hổng thực trong 68
dự án C nguồn mở. 58.2\% các lỗ hổng được phân loại nằm ở những kiểu mà Rust an toàn gần như
đã làm không thể xảy ra~\cite{meneely2025justuserust}. Tỉ lệ đó là trường hợp tốt nhất: nó giả
định mã mới ở lại trong Rust an toàn. Một kiểu được tính là đã ngăn nếu ánh xạ nói vậy, dù
không trình biên dịch nào đã từ chối một chương trình. RQ2 hỏi cùng câu hỏi đó với các chương
trình được biên dịch và chạy. Đó là cách nó có thể nói khi nào các kiểu còn lại bị bắt, và tốn
gì để bắt chúng.
